\textbf{Eligibility.} A \emph{core study} (i) analyses or intervenes on an internal representation
(\S\ref{sec:prelim}) of (ii) an eligible medical vision or vision--language foundation model, or an
encoder used to assemble one (\S\ref{sec:models}), and (iii) reports a measurement. Three decisions are
kept apart in the eligibility coding. \emph{Foundation-model status} is the definition of
\S\ref{sec:models}, operationalised for sampling as pre-training on at least $10^{5}$ images or image--text
pairs, $10^{3}$ whole slides or $10^{4}$ volumetric scans from more than one anatomical site or finding
class, and reuse across at least three reported tasks. The three thresholds sit below the smallest
pre-training corpus of each kind in Table~\ref{tab:models}, so that they exclude task-specific models
rather than small foundation models. They compare information units that are not equivalent, since an
image, a slide and a volume differ by orders of magnitude in content, so they are sampling boundaries
rather than part of the definition, and the conclusions do not depend on them: excluding the six
borderline studies changes none of the statements of Table~\ref{tab:statements}; excluding also the
seventeen studies admitted by the volumetric threshold empties the brain-MRI row of the site finding
(F3c) and narrows the modality span of two others (F2 and F8a$^{\prime\prime}$) without reversing
either. Medical adaptations of general foundation models qualify, as do their component encoders.
\emph{Representation access} is whether the study could reach embeddings or activations at all. It is a
requirement of this review rather than a property of the model, and it biases the corpus towards openly
released weights, since API-only models enter only through embedding-level analyses and are marked
borderline. \emph{Claim eligibility} is whether the study makes a representation-level claim about such
a model. One family carries a stricter version of it: a geometry study enters only
if it connects its statistic to decodability, retrieval, transfer or shortcut behaviour
(\S\ref{sec:geometry}). The structural family is therefore reviewed as \emph{task-linked geometry}, and
purely descriptive or negative geometry results are under-represented here relative to the literature. A
study enters if at least one analysed model passes the first two decisions and the study passes the
third; ineligible comparators are named in the evidence tables and support no finding alone.
General-domain methods, datasets, benchmarks, model papers and surveys are contextual. Output-only and
input-attribution-only studies are excluded unless the attribution tests a representation-level claim;
so are standalone clinical language models, non-imaging modalities and task-specific networks. Model and
training-method papers in which the representation measurement is only instrumental, an attention map
shown as qualitative evidence or an audit that selects heads for fine-tuning, are contextual, and the
size of the tracing family depends on this rule (\S\ref{sec:tracing}). Publication status is not a
criterion: preprints, theses and repository deposits with a retrievable full text are coded like refereed
papers and counted as preprints. Frozen-feature benchmarks are admitted only when they report a
representation-level property beyond task accuracy; accuracy leaderboards, and four parameter-editing
benchmarks that report only output-level reliability, locality and generality, are contextual.

\textbf{Search and screening.} The literature was searched between 15 and 18 August 2026 in Semantic
Scholar, OpenAlex and arXiv, for work published from January 2018 to the cutoff of 16 August 2026, and
MEDLINE (PubMed) was searched on 21 August 2026. Each query combined the vocabulary of one method family
(probing and concept decoding, concept bottlenecks, dictionaries and dissection, attention and
attribution, steering, patching and editing, embedding geometry, adaptation, evaluation) with
medical-domain and foundation-model terms, in twelve query groups. The database searches were followed
by backward and forward citation chaining~\cite{wohlin2014guidelines} in five rounds, seeded from
\Nmodelseeds{} model papers, among them every model of Table~\ref{tab:models}, and from every included
study; chaining contributed \Nstagechain{} of the \Ncore{} core studies, whereas the MEDLINE search,
which indexes few of the preprints, retrieved \pmRecall{} of them and added none. Records were
de-duplicated across sources by identifier and normalised title. In all, about 13,600 records were
screened on title and abstract, \Nretained{} were retained for full-text assessment, and \Ncore{} met
the criteria above as core studies; four of these, whose full text is closed access, were coded from the
abstract and the openly available figures, tables and supplementary material. The article cites
\Ncited{} references in all: the core studies and \Ncontext{} contextual references, drawn from the
retained records and from \Ncitedhand{} documents added by hand (regulatory guidance and technical
reports that no database returns). Screening and coding were done by one reviewer, so no inter-rater
agreement is reported; borderline eligibility decisions were recorded as such, and
the query strings, the counts of every stage and the decision recorded for every screened record are in
the online supplement.

\begin{figure*}[t]
\centering
\includegraphics[width=\linewidth]{fig2_corpus.pdf}
\caption{The reviewed corpus. (a) Core studies per year of first appearance, by primary method family
(\Ncore{} studies; cutoff 16 August 2026, so 2026 is partial). (b) The \Nclaims{} records by family
and locus: the vision encoder is split by sublocus, Succ.\ loci is a probe repeated at successive loci,
the intervention family is split into interventions inside the forward pass and post-hoc operations,
and hatched cells contain no record; the interventional records are the twenty-one of the
causal-intervention family and one decoding record with a counterfactual end-point. (c) Records by
modality (Other: echocardiography, PET, cytology, endoscopy and mixed-modality retrieval), with records
spanning modalities counted under each. The structural family is a sample of task-linked geometry
studies, marked by a dashed outline in (a) and (c), and its size is not comparable
with the other families. Proportions describe the assembled corpus, not the literature.}
\label{fig:corpus}
\end{figure*}

\textbf{Coding.} Two units are used throughout. A \emph{record} is one result of one study: one study,
one method family, one evidence type and one locus, so that outcomes sharing all four form a single
record, which may still span several concepts, models, datasets or end-points. A \emph{statement} is a
cross-study proposition of \S\ref{sec:findings}, and the records that support it are listed under it.
Each record was coded for its family and sub-group, locus, concept provenance, evidence type, models,
modality, whether it compares a medical with a general model on the same data, whether it evaluates more
than one model or dataset, and, for interventional claims, the control used (\S\ref{sec:taxonomy}).
Every count in this article is a count of records or of studies in this corpus.

\textbf{Where the evidence sits.} The corpus holds \Ncore{} core studies and \Nclaims{} records;
\Nrecent{} of the studies (\Pctrecent\%) appeared in 2025--2026, and 2026 is a partial year
(Fig.~\ref{fig:corpus}a). Fig.~\ref{fig:corpus} shows the shape of the corpus. Observational
claims at the vision encoder and in the projected joint space dominate, pathology and chest radiography
supply nearly two thirds of the records, and the few interventional claims divide between the decoder
residual stream and post-hoc operations on stored embeddings. Only one record repeats a probe at
successive loci~\cite{zhu2026lost} and one is taken at the text encoder~\cite{tayebiarasteh2026self};
none is taken at the connector output alone or at the pre-readout representation (\S\ref{sec:future}).
The synthesis is descriptive: no effect is pooled, no certainty is graded, and an empty cell in
Fig.~\ref{fig:corpus}b means that no study was found, not that none exists. Study quality is described
by the controls and reporting practices the studies themselves state (\S\ref{sec:findings}); the
risk-of-bias instruments of prediction modelling (PROBAST and
PROBAST+AI~\cite{wolff2019probast,moons2025probast}, TRIPOD+AI~\cite{collins2024tripod},
QUADAS-2~\cite{whiting2011quadas}) are not designed for representation-level claims.
